\section{실제 계산과 독립 검증}
\subsection{동결 입력, 실행과 정밀도}
PHYS18의 최종 commit과 연구 묶음을 기준으로 시작했다.
채택된 곡률 계약의 23개 입력은 경로·바이트 수·SHA-256이 모두 일치했다.
PHYS15--18의 기존 source box, root·carry·inverse 및 response 포함을
그 원본 의미와 identity에 결박해 재사용했다.

\texttt{src/curvature\_application.py}는 정확한 Fraction 연산에
256-bit dyadic 외향 반올림을 결합한다.
순수 jet의 곱과 역수에는 보통 2차 도함수 계수를 사용했다.
첫 half·Full·carried 둘째 half,
직접·상관 endpoint forcing,
35개 photon slot, 일곱 sector, number·energy ledger와 온도 quotient를
한 번의 실제 계산에서 출력했다.
실행 전에 명령·code hash·환경·시간 예산을 기록했다.

\begin{table}[H]\centering\small
\caption{완료한 계산의 범위와 실제 비용.}
\begin{tabularx}{\textwidth}{lXl}
\toprule
실행 & 계산 범위 & 경과\\
\midrule
application run01 & 동결 원 입력의 순수 곡률 포함 생성 & wrapper 1.569 s\\
application 내부 & 동일 실행의 실제 계산 구간 & 1.407 s\\
geometry readout02 & 정확한 유리수 Hessian·보간·계수 예산 & 0.049 s\\
독립 Fraction run03 & 별도 산술의 실제 포함·readout 검증 & wrapper 1.870 s\\
Wolfram run01 & 실제 xAct와 symbolic 순수 항등식 & API 13.222 s\\
Wolfram 내부 & 같은 요청의 kernel 대수 계산 & 2.000 s\\
Wolfram run02 & decimal 입력을 exact rational로 한 Hessian 방향 검산 & API 6.088 s\\
\bottomrule
\end{tabularx}
\end{table}
Python 환경은 3.12.14, Linux x86-64다.
main wrapper의 최대 RSS는 55,992 KiB, 계산 내부 기록은 47,544 KiB다.
geometry의 readout01은 계산 결과를 먼저 생성했고,
readout02에서는 점별 혼합 미분과 corner 적분 계수의 입력 필드를 분리했다.
두 SUMMARY의 정확한 endpoint는 전부 같았다.
원 실행과 두 readout 기록을 각각 보존했다.

새 atomic source 호출, reaction fit, nonlinear root solve 및 history replay의 횟수는
각각 0이다. 이번 실행의 대상은 이미 검증된 원 입력에서 새 순수 도함수와
유한 응답 포함을 계산하는 것이었다.

\subsection{Wolfram과 xAct의 실제 사용}
연결된 Wolfram 15.0.1 Linux kernel에서 공식
xAct 1.3.0 archive를 내려받아 hash를 대조하고
xTensor 1.3.0·xPerm 1.2.4를 실제 로드했다.
tensor canonicalization은 \texttt{\$xpermQ=False}인 Wolfram Language group backend다.
패키지 초기화의 external xPerm 메시지와 실제 선택한 backend를 구분해 기록했다.

첫 요청의 6개 tensor 항등식, 13개 scalar 항등식,
14개 matrix component가 모두 literal 0을 반환했고,
Peano kernel의 부호·질량·weight에 관한 4개 부등식은 모두 True였다.
검산은 implicit 계수 2,
carried $Q_2''$,
Full birth 계수,
대칭 Hessian 차이,
원 순서의 일곱 photon sector,
유한 온도 quotient와 보간 kernel을 포함했다.
matrix sector 검산에는 임의의 비가환 2-slot remap을 사용했다.
35개 실제 photon slot의 수치 포함은 별도의 application과 독립 심사에서 확인한다.

첫 요청에는 임시 symbol 이름에 대한
\texttt{Symbol::undefined2} 진단이 있었다.
제출된 Module의 local 이름과 대입을 확인했고,
지정한 residual은 모두 평가된 0이었다.
원 경고·초기화 메시지·request·response를 보존했다.
두 번째 Hessian 요청은 경고 없이 완료됐으며,
별도의 두 방향 수축 항등식도 정확히 0이었다.
decimal 표시를 exact rational로 사용한 별도 포함은
main의 원 유리수 포함보다 약간 넓지만 동일한 엄격한 부호를 준다.

\subsection{독립 심사와 채택}
독립 원식 유도는 $C^2$ 계약에서 C3--C20의 부호·계수를 확인했다.
별도 Fraction interval class, 이항계수 pure jet, Gaussian elimination으로
실제 producer의 산술을 import하지 않고 포함을 재계산했다.
최종 run03에서 362개 검사,
265개 비교 묶음의 2,057개 스칼라 포함,
3,174개 외향 decimal 표시,
23개 계약 입력이 모두 통과했다.
같은 Full inverse, 35개 carried photon slot과 일곱 sector,
endpoint·온도 quotient 및 유한 응답을 독립적으로 확인했다.
number·energy ledger의 별도 수치 잔차는 모두 0을 포함했으며,
정확한 보존식의 증명은 원식 유도에 연결했다.

독립 검산 wrapper의 경과는 1.870 s, 최대 RSS는 52,448 KiB다.
앞선 두 짧은 시도는 검사기의 정규화 비교와 변수 shadowing을 고친 뒤 종료했고,
세 실행의 snapshot·실패 기록을 모두 보존했다.
수정 대상은 독립 검사기였으며 scientific producer와 원 입력·결과는 그대로다.
최종 독립 판정은 \texttt{PASS\_SCOPED}다.
Root는 검토된 source·result·readout hash를 다시 대조한 뒤,
현재 고정 계열에서 \texttt{PROMOTE\_SCOPED}로 채택했다.

원 유도 문서는
\texttt{theory/PURE\_CURVATURE\_INDEPENDENT\_DERIVATION\_KO.md}에
보다 세부적인 C3--C20 대응과 모든 부호·극한 검토를 보존했다.
HH source가 항등적으로 0인 경우의 $\lambda$ 응답,
birth가 0인 경우의 $b$ 응답,
투명 opacity 극한,
같은 endpoint 계열에서의 결함 소멸,
bilinear 함수와 corner에서의 정확한 재현을 원식으로 점검했다.
이 극한들은 실제 source의 추가 실행으로 보고하지 않았다.

\subsection{문헌 대조가 지지하는 부분}
SciSpace의 완전한 질문 1회와 Exa의 두 검색(각 6개 후보)에서
관련 문헌을 조사하고, 세 원전의 실제 해당 본문을 대조했다.
Exa에서 RWTH PDF의 본문 회수가 비어 나온 한 사례는
별도 web 원문 열기로 확인을 마쳤다.
상세 요청·후보·접근 상태와 읽은 위치는 문헌 provenance에 있다.

일반 implicit 2차 미분 \cite{SLN15},
검증된 interval inverse의 family 양화 \cite{R99},
순수 편미분만을 사용하는 tensor-product 보간 경계 \cite{MR09}는
기존 수학 구조다.
이번의 구체적인 연구 진전은 원 gas--photon carry,
HH와 birth의 전역 parameter identity,
상관된 endpoint 잔차와 유한 온도 quotient를 결합해
실제 Kelvin 단위의 부호와 오차 포함을 닫은 데 있다.

\section{다음 루프의 물리 질문과 인계}
\subsection{PHYS20: 실제 소비 구간의 HH source 응답}
다음 목표는 실제로 이름과 원 입력을 지정한
rei\_bianchi 소비 구간 하나에서
HH의 절대 열 변화율과 전자 생성률을 구하는 것이다.
첫 행동은 최신 채택 entry에서 해당 구간의 source,
COMMON, 상태와 단위, density/frame, clock,
geometry/tilt, closure와 process owner를 좁게 회수해 결박하는 것이다.
현재 consumer 결박은 이번 루프에서 확인하지 않았으므로
\texttt{unverified}로 인계한다.

원 readout과 해당 소비자의 HH stoichiometry가 채택되면
\[
\dot g_{\rm HH}=\nu_Hq_{\rm HH},\qquad
\dot T_{\rm HH}
=\frac{C_Te_w^T-T\ell_e}{\Pi}\nu_Hq_{\rm HH},\qquad
\dot e_{\rm HH}=\ell_e\nu_Hq_{\rm HH}
\tag{19.31}
\]
를 외향 포함한다.
여기서 $q_{\rm HH}$는 실제 HH 활성 계수까지 포함한 이벤트율이다.
원 FT03/LCS의 $\nu_H=(1,0,0,-\chi_H)^T$가 admitted되면
\[
\boxed{\dot T_{\rm HH}
=-\frac{C_T\chi_H+T}{\Pi}q_{\rm HH},\qquad
\dot e_{\rm HH}=q_{\rm HH},\qquad
(\dot n_e)_{\rm HH}=n_Hq_{\rm HH}.}
\tag{19.32}
\]
단위는 각각 K/s, per H/s, 전자수밀도/s다.
consumer가 $c^{-1}\partial_t$ 정규화를 사용한다면 원 $c$와 clock을 먼저 변환한다.
첫 표에는 같은 상태의 다른 thermal/electron source, signed total,
HH의 절대 기여와 지배 입력을 함께 담는다.
source ownership을 확인해 중복 합산을 피한다.

이 목표는 기존 \texttt{HH\_DOMAIN\_SCREEN}에서
\texttt{HH\_MATERIALITY}로 이어진다.
local source 응답에 필요한 적용성 결박과
PHYS19의 전체 endpoint 인증을 재사용하는 적용성 결박을 각각 그 claim에 맞게 확인한다.
finite interval·matched history는 기존 root/carry 및 RUX11 조건이 결박되면 이어간다.
준비된 인계 계약은 실제 다음 계산을 수행한 결과로 기록하지 않는다.

\subsection{프로젝트 상태와 재현 진입점}
이번 채택은 \texttt{HH\_DEFECT\_PARAMETER\_CURVATURE}의 범위에 한정한다.
기존 RUX 전역 DAG와
\texttt{HH\_STRICT\_NATIVE}, \texttt{HH\_SINGLET}의 독립 질문을 유지한다.
이 두 lane의 원 microscopic/channel/spin 조건은 각자 원 입력에 연결한다.
PHYS20 소비 source 응답과의 완료 의존성을 새로 만들지 않는다.

제공한 묶음에는 보고서 소스·수치 JSON·실행 코드·원 request/response,
입력 identity와 PHYS18 dependency archive, 독립 심사와
다음 consumer 계약이 들어 있다.
최상위 \texttt{REPRODUCE.md}는 이미 완료된 근거를 확인하는 경로와
필요할 때 같은 새 계산만 재현하는 명령을 제시한다.
문서·수치 포함·원 source/root 근거의 상태는
\texttt{review/DECISION.json}과 \texttt{state/RESEARCH\_STATE.json}에 분리해 결박했다.
